\chapter{接收通路}\label{ch:rx}

接收核心 \code{afdx_rx_core} 为 A、B 两路网络各配置一套接收 MAC 和帧头解析器，两网共用同一接收引擎与交付模块。MAC 在接收过程中逐字节校验 CRC，解析器同步检查帧头并将载荷预写入本网的环形缓冲。对于通过解析器检查的帧，接收引擎根据帧描述依次执行 VL 查表、完整性检查和冗余管理，再通知解析器提交数据或回退写指针；已提交的消息最终由交付模块逐字节送往应用口。

\section{接收核心与配置表}\label{sec:rxtables}

\figN[接收核心框图]{f11_rx}{接收核心及配置、状态和载荷存储之间的连接}{fig:rx}

接收侧的片上表见表~\ref{tab:rxtables}。\code{rx_map} 以 16 位 VLID 直接寻址，任意 4096 个 VLID 均可映射到接收 VL 索引 \code{idx}；映射表定义为 65536$\times$13 的存储阵列，读出数据直接锁存至一级寄存器，查表延迟固定为 2 拍。

\begin{table}\centering\small
\caption{接收侧配置与状态表}\label{tab:rxtables}
\begin{tabularx}{\linewidth}{>{\raggedright\arraybackslash}p{2.0cm}>{\raggedleft\arraybackslash}p{1.9cm}>{\raggedright\arraybackslash}p{1.9cm}>{\raggedright\arraybackslash}p{1.9cm}L}
\toprule
表 & 深度$\times$位宽 & 写口 & 读口 & 内容 \\
\midrule
\code{rx_map} & 65536$\times$13 & 配置总线 & \code{rx_engine} & \code{valid}[12]、\code{idx}[11:0] \\
\code{rx_cfg} & 4096$\times$32 & 配置总线 & \code{rx_engine} & \code{en}、\code{net_a}、\code{net_b}、\code{ic_en}、\code{rm_en}、\code{skew_max}[16:5]、\code{max_payload}[27:17] \\
\code{rx_state} & 4096$\times$44 & \code{rx_engine} & \code{rx_engine} & 每网 PSN 及有效位；\code{rm_last}、\code{rm_time} 及有效位 \\
\code{rx_ring_a/b} & 4096$\times$8 & 本网 \code{rx_parse} & \code{rx_deliver} & 每网一个环形缓冲，按到达顺序存载荷 \\
\code{rx_dq} & 32 项寄存器 & \code{rx_engine} & \code{rx_deliver} & 待交付消息的来源网络、起始地址、长度、\code{idx}、UDP 端口 \\
\bottomrule
\end{tabularx}
\end{table}

\code{rx_cfg} 中的 \code{skew_max} 以 \code{now_c} 的计时单位（32.768~µs）计。\code{afdx_rx_core} 例化 A、B 两路接收 MAC 与解析器、共用接收引擎、交付模块及 5 张表，其中 \code{rx_map}、\code{rx_cfg} 的写口由核心统一管控。复位后，核心用 65536 拍逐项清零 \code{rx_map}（前 4096 拍同步清零 \code{rx_cfg}），引擎另行清零 \code{rx_state}；全部完成后置位 \code{init_done}。运行期间仅在 \code{cfg_ok} 为 1 时锁存配置写请求并于下一拍写入配置表。核心将各模块的统计脉冲按位相或，寄存一拍后输出。

\section{接收 MAC}\label{sec:macrx}

\code{afdx_mac_rx} 把 GMII 输入寄存一拍后，在 \code{rx_valid} 为 1 的每一拍处理一个字节，状态机见图~\ref{fig:macfsm}(b)。
\begin{itemize}
  \item \textbf{帧起始定界（SFD 检测）：}前导字节 \code{0x55} 的数量不限，检测到 \code{0xD5} 后以紧随的下一字节作为帧首字节。若收到其他值，或在检测到 \code{0xD5} 之前 \code{rx_valid} 已回到 0，则判为前导码错误。复位后若 \code{rx_valid} 已为 1，MAC 先等待其回到 0，再从下一帧开始接收。
  \item \textbf{剥离 FCS：}每接收一字节均更新 CRC 和帧长；输出比输入晚 4 字节，收到第 $k+4$ 字节时输出第 $k$ 字节，最后 4 字节（FCS）不送后级。
  \item \textbf{帧结束：}寄存后的 \code{rx_valid} 首次为 0 时，MAC 判定帧接收结束，输出 \code{f_end} 和 \code{f_err}。错误判定依次检查前导码错误、帧长不在 64--1518 范围内、CRC 余式不等于 \code{0xDEBB20E3}，均无异常则判为正确帧。MAC 向后级最多输出 1514 字节；若帧长超限，则按帧长错误结束本帧。
  \item \textbf{帧间：}\code{f_end} 之后立即回到空闲，可分辨间隔为 1 个低电平周期的两帧。
\end{itemize}
以 GMII 上 FCS 末字节到达时刻为第 0 拍：第 1 拍该字节锁存至输入寄存器，第 2 拍寄存后的 \code{rx_valid} 拉低，第 3 拍输出 \code{f_end}、\code{f_err} 和 \code{f_len}。剥离 FCS 后的最后一个有效字节在第 2 拍从 \code{m_data} 输出，比第 3 拍的 \code{f_end} 早 1 拍，解析器据此完成校验和累加。

\section{帧头解析与预写入缓存}

\code{afdx_rx_parse} 由字节位置计数器 \code{position} 驱动，随输入字节流式检查帧头；另设两个辅助状态机分别负责长度计算与帧尾分类，结构见图~\ref{fig:parse}。

\subsection{流式检查与长度计算}

\figN[帧头解析]{f23_parse}{帧头解析：(a) 按字节位置检查并预写入载荷；(b) 收到 UDP 长度后的两拍长度计算；(c) 帧尾的两拍分类}{fig:parse}

帧头必须满足以下条件，任一条件不满足均判为“帧头错”：
\begin{itemize}
  \item 第 0--3 字节为 \code{03 00 00 00}；第 12--13 字节为 \code{08 00}；第 14 字节为 \code{0x45}；第 23 字节为 \code{0x11}；
  \item IP 不分片（MF 为 0、片偏移为 0）；IP 头 10 个 16 位字的反码和为 \code{0xFFFF}；
  \item 记 $n$ 为 UDP 长度减 8，要求 $1\le n\le 1471$，IP 总长等于 $n+28$，且含 FCS 的帧长等于 $47+\max(n,17)$；
  \item UDP 校验和非 0 时，伪头、UDP 头与 $n$ 字节数据的反码和为 \code{0xFFFF}。
\end{itemize}
补零字节的值不检查。解析器同时记录 VLID、源/目的端口、$n$ 和 SN。

在第 39 字节收到 UDP 长度低字节后，\code{length_phase} 控制长度计算分两拍完成：先在 LEN\_SIZES 步计算 $n$、检查长度是否合法，并求出期望 IP 总长；再在 LEN\_LIMITS 步计算载荷末位置、SN 位置、期望帧长和环形缓冲所需空间。第 42 字节到来时，上述结果均已就绪。

IP 校验和与 UDP 校验和分别使用 20 位和 26 位累加器，每两字节累加一次：前者累加 IP 头的 10 个 16 位字，后者累加 UDP 伪头、UDP 头及载荷。进位折叠采用两级寄存器流水，与累加过程并行；到帧尾后第 1 拍（FIN\_SETTLE），UDP 校验和结果最终稳定。

\subsection{载荷预写入与提交}

从第 42 字节开始，解析器将载荷边接收边写入本网环形缓冲，待最终判定后再决定是否提交。环形缓冲采用三个指针（图~\ref{fig:ring}）：读指针 \code{rd} 由交付模块维护；已提交写指针 \code{wr_commit} 和预写入指针 \code{wr_trial} 由本网解析器维护。

计算剩余空间时以 \code{wr_commit} 为准，不计入正在接收的帧。第 42 字节到来时 $n$ 已知，若剩余空间小于 $n$，解析器便不写入本帧载荷，并将其记为溢出。

\figN[接收环形缓冲的三个指针]{f12_ring}{接收环形缓冲的三个指针。帧通过全部检查后，解析器将 \texttt{wr\_commit} 更新为 \texttt{wr\_trial}；否则将 \texttt{wr\_trial} 回退至 \texttt{wr\_commit}}{fig:ring}

\subsection{帧尾分类与描述交接}

每次 \code{f_end} 有效时，解析器均对该帧进行一次分类。\code{finish_phase} 在 \code{f_end} 后第 2 拍进入 FIN\_REPORT 并作出判定，类别脉冲和帧描述均经寄存后输出，在第 3 拍有效。分类时，若 \code{f_err} 表示 FCS 错误或帧长错误（含前导码错误），则直接采用该结果；否则依次判定帧头错误、溢出或正确。

解析器在帧尾完成错误帧的分类与计数，并将预写入指针 \code{wr_trial} 回退至已提交写指针 \code{wr_commit}；正确帧则生成帧描述字交由接收引擎处理（表~\ref{tab:ifdesc}）。两网在同一拍结束的正确帧，其帧描述也在同一拍给出。

同一网相邻两帧的 \code{f_end} 至少相隔 73 拍（最短帧 64 字节加前导码 8 字节和 1 拍间隔）。引擎处理一份描述至多需要 13 拍，即使两网各有一份描述，也只需 26 拍即可处理完毕，短于同一网相邻描述的到达间隔。若数据帧的首个前导字节在清表完成前到达，解析器将忽略整帧，不生成帧描述，也不计入统计。

\section{完整性检查与冗余管理}

\code{afdx_rx_engine} 是 \code{rx_state} 的唯一读写控制模块，一次处理一份帧描述或执行一步冗余记录过期清理，两网描述同拍到达时先处理 A 网。

\subsection{帧描述处理流程}

帧描述的处理流程见图~\ref{fig:decide}。引擎先查询 \code{rx_map}，再读取 \code{rx_cfg} 和 \code{rx_state}，随后依次执行配置检查、完整性检查、冗余管理和交付入队。任一步未通过，均按对应原因计数，并通知解析器回退预写入指针。

\figN[接收引擎对一份帧描述的判定流程]{f13_decide}{接收引擎对一份帧描述的配置检查、完整性检查、冗余管理与交付判定}{fig:decide}

引擎有 INIT、IDLE、FRAME、AGE 四个操作态，FRAME 按 \code{phase} 0--11 逐拍推进，每拍至多一级比较，见图~\ref{fig:rxeng} 和表~\ref{tab:rxphase}。对于未通过检查的帧，引擎提前跳转至写回或报告步骤，因此一份描述的处理时间为 5--13 拍（含 IDLE 取描述的 1 拍）。图中的 \code{desc} 为 \code{desc_a_valid} 或 \code{desc_b_valid}；AGE 的任一步收到新描述时均返回 IDLE。

\figN[接收引擎状态机]{f24_rxeng}{接收引擎状态机：FRAME 处理帧描述，AGE 在空闲时检查冗余记录}{fig:rxeng}

\begin{table}\centering\small
\caption{接收引擎 FRAME 各步}\label{tab:rxphase}
\begin{tabularx}{\linewidth}{>{\centering\arraybackslash}p{1.2cm}L}
\toprule
\code{phase} & 动作 \\
\midrule
（IDLE） & 取走帧描述（\code{take}），以 VLID 读 \code{rx_map} \\
0--1 & 等待同步读出，保存映射 \\
2 & 映射无效 $\rightarrow$ 11（未知 VL）；否则以 \code{idx} 读 \code{rx_cfg} 和 \code{rx_state} \\
3--4 & 等待读出，保存配置和状态 \\
5 & VL 或本网未启用、$n>$\code{max_payload} $\rightarrow$ 11（配置不允许）；同时计算 PSN 的一步后继序号、到达时间差、SN 与 \code{rm_last} 的原始距离 \\
6 & 时间差取绝对值，计算 PSN 的两步后继序号，原始距离折算到 1--255 环上 \\
7 & 完整性检查，更新本网 PSN；不通过 $\rightarrow$ 10 \\
8 & 冗余管理判定 \code{deliver_allowed} \\
9 & 执行冗余去重判定；若 \code{rx_dq} 已满则记为溢出，否则写入交付队列并更新 \code{rm_last}、\code{rm_time} \\
10 & 状态有变化则写回 \code{rx_state} \\
11 & 给出 \code{resolve} 与 \code{commit}，按处理结果产生一个统计脉冲，返回 IDLE \\
\bottomrule
\end{tabularx}
\end{table}

\FloatBarrier
\subsection{配置与完整性判定}

\code{rx_map} 映射无效时，引擎将该帧计为“未知 VL”；VL 未启用、本网未启用或 $n$ 大于 \code{max_payload} 时，计为“配置不允许”。

完整性检查（\code{ic_en}$=1$ 时）每网独立进行。记本网上一帧的序号为 PSN，满足以下任一条件即可通过检查：清表后本网首次收到该 VL；SN 为 0（对端复位）；SN 为 PSN 在 1--255 环上递增 1 步或 2 步的后继序号，后者允许丢一帧。PSN 为 0 时，后继序号分别为 1 和 2。

启用完整性检查时，通过映射与配置检查的帧均以当前 SN 更新本网 PSN，包括序号检查失败的帧；后续序号连续时可恢复同步。未通过序号检查的帧计为“完整性检查丢弃”。

引擎先对每份副本独立进行完整性检查，通过后再进行冗余管理，符合 ARINC 664 P7 的处理顺序。

\subsection{冗余窗口与过期清理}

冗余管理（\code{rm_en}$=1$ 时）按处理顺序交付首个有效副本；两网描述同拍待处理时先处理 A 网。副本到达间隔在 \code{skew_max} 内时，同一消息只交付一次，判据如下：
\begin{itemize}
  \item 尚无记录，或到达时刻与 \code{rm_time} 之差超过 \code{skew_max}：交付；当两份副本的处理顺序与到达顺序相反时（如引擎忙期间 B 网先完成解析、A 网后完成解析，引擎转入空闲时按 A 网优先仲裁先处理了 A 网），对 16 位时间差取绝对值。
  \item 否则 SN 等于 \code{rm_last}：另一网的重复帧，丢弃并计“冗余管理丢弃”。
  \item 否则 SN 为 0 或 \code{rm_last} 为 0：交付。
  \item 否则在 1--255 环上，SN 领先 \code{rm_last} 1--127 步则交付，其余丢弃。
\end{itemize}
对于判定可交付的帧，若 \code{rx_dq} 已满，则计为“溢出”并回退预写入指针；否则提交环形缓冲中的载荷，将帧描述写入 \code{rx_dq}，并更新 \code{rm_last} 和 \code{rm_time}。两网描述同拍到达时，从描述有效到引擎完成判定，后处理的一份也不超过 27 拍（设计约束取 64 拍）。

\code{now_c} 每变化一次，引擎在空闲时通过 AGE 状态的 4 个步骤检查下一条 VL；有帧描述到来时立即返回 IDLE 处理帧，之后重试该项过期清理检查。判据为：若记录已超过 \code{0x4000} 个计时单位（约 0.537~s），清除 \code{rm_time_valid}。

帧描述处理优先于过期清理，清理被打断后会重新检查同一项。4096 项约需 0.134~s 检查一遍，记录在约 0.671~s 内失效，早于 16 位时间差的半圈时间 1.074~s。

\section{交付}\label{sec:rxdeliver}

\code{afdx_rx_deliver} 按入队顺序从 \code{rx_dq} 取出帧描述，再从对应环形缓冲每拍读出并输出一个字节，同时释放已读空间，结构见图~\ref{fig:deliver}。

\figN[交付模块]{f25_deliver}{交付模块的队头选择、环形缓冲读取与字节输出流水}{fig:deliver}

图~\ref{fig:deliver} 中 C1--C6 依次为组内选择寄存、队头寄存、读请求寄存、同步 RAM 读出、两网数据寄存和输出寄存六个采样边界。队头选择状态机按 \code{HEAD_IDLE}（组内 8 选 1，C1）、\code{HEAD_SELECT}（组间 4 选 1 锁存至 \code{dq_data} 并给出 \code{dq_pop}，C2）和 \code{HEAD_RELEASE}（等待出队握手完成）推进；队头进入 \code{dq_data} 后，读控制在下一拍发起首次读取，随后每拍将地址加 1，直到 \code{left} 减为 0 时停止。

读请求、同步 RAM、数据寄存三级流水由 \code{valid_pipe}、\code{last_pipe} 对齐，第一个 \code{d_valid} 在 \code{dq_count} 变为非零后的第 6 拍有效。\code{d_last} 有效的当拍，模块重新允许取队头。

一条消息的 $n$ 个字节在连续 $n$ 拍内送出，\code{d_last} 与末字节同拍；VL 索引、端口和来源网络在整条消息期间不变。交付口不带反压。

线路空闲时，从帧尾到首个 \code{d_valid} 的典型延迟为 24 拍，设计上界为 200 拍（图~\ref{fig:pipe}）。若前一消息尚未交付完，当前消息需等待；取队头与读流水共约 6 拍，可在前一消息末字节之后的 16 拍内开始交付。32 项描述队列用于暂存长载荷交付期间到达的短帧描述（第~\ref{sec:budget}~节）。

\FloatBarrier
\Needspace{22\baselineskip}
\section{统计分类}\label{sec:rxclass}

对于清表完成后开始接收的每一帧（一段 \code{rx_valid} 连续为 1 的字节流），解析器均在帧尾离开线路后的 32 拍内，在对应网络上输出且仅输出一个类别脉冲。对于分类为正确的帧，接收引擎还输出且仅输出一个处理结果脉冲；该脉冲在引擎完成判定后的 8 拍内有效，且不晚于帧尾后的 150 拍。每网的 \code{ev} 位定义见表~\ref{tab:rxev}。

\begin{table}\centering\small
\caption{接收统计脉冲 \texttt{ev\_a}、\texttt{ev\_b}}\label{tab:rxev}
\begin{tabularx}{\linewidth}{>{\centering\arraybackslash}p{1.0cm}>{\raggedright\arraybackslash}p{3.0cm}>{\raggedright\arraybackslash}p{2.2cm}L}
\toprule
位 & 名称 & 产生模块 & 含义 \\
\midrule
0 & 正确帧 & \code{rx_parse} & 通过 MAC 和帧头检查 \\
1 & FCS 错误 & \code{rx_parse} & CRC 校验不符 \\
2 & 帧长错误 & \code{rx_parse} & 帧长越界、超长或前导码错误 \\
3 & 帧头错误 & \code{rx_parse} & 固定字段、IP/UDP 长度或校验和错误 \\
4 & 未知 VL & \code{rx_engine} & \code{rx_map} 无效 \\
5 & 配置不允许 & \code{rx_engine} & VL 或本网未启用，或超过 \code{max_payload} \\
6 & 完整性检查丢弃 & \code{rx_engine} & SN 不连续 \\
7 & 溢出 & 两处相或 & 环形缓冲剩余空间不足，或 \code{rx_dq} 已满 \\
\bottomrule
\end{tabularx}
\end{table}

位 7 由两处溢出脉冲相或产生：解析器报告环形缓冲剩余空间不足，引擎报告 \code{rx_dq} 已满。同一网上，相邻帧尾至少相隔 73 拍，引擎判定至多 27 拍，解析器的第 $k$ 帧类别脉冲与引擎的第 $k-1$ 帧结果脉冲分时到达，可相或合并。

两网合计另有 \code{ev_rm_drop}（冗余管理丢弃）和 \code{ev_deliver}（交付）。这些脉冲与第~\ref{ch:cfg}~章的计数器一一对应，记录每类接收结果。
